\chapter{Marco Teórico}
(En este apartado el estudiante debe presentar en forma breve los temas que se abordarán en el marco teórico. El marco teórico se desprende de las palabras clave y del titulo. Debe ser totalmente acotado.
Ésta sección describe conceptos relacionados a procesamiento de trayectorias y detección de trayectorias anómalas .

\section{Primer subtitulo}
Aquí debe llenar tu contenido

\begin{algorithm}[H]
\SetAlgoLined
\KwData{$G=(X,U)$ such that $G^{tc}$ is an order.}
\KwResult{$G’=(X,V)$ with $V\subseteq U$ such that $G’^{tc}$ is an interval order.}
initialization\;
\nl \Begin{
\nl\While{not at end of this document}{
	\nl read current\;
	\nl \eIf{understand}{
	\nl go to next section\;
	\nl current section becomes this one\;
	}{
	\nl go back to the beginning of current section\;
	}
 }
 }
\caption{Ejemplo de como colocar código en código latex}
\label{IR2}
\end{algorithm}

Ejemplo para colocar citaciones: \cite{sillito2008semi}

Es posible utilizar abreviaturas dentro del contenido, abreviaturas como:
%\acronum{GPS}, \ac{INEI}, \ac{t-SNE}, \ac{SOM}, \ac{IBM}.

Para la utilización de ecuaciones matemáticas se puede utilizar el simbolo \$\$ dos veces tanto al inicio y al final, y para otros casos utilizar .

$$ E = mc^2 $$

Ecuaciones:
\begin{align*}
  1 + 2 &= 3\\
  1 &= 3 - 2
\end{align*}

Matrices: 
\[\begin{bmatrix}
1 & 0\\
0 & 1
\end{bmatrix}\]

otros:
\begin{equation}
(\frac{1}{\sqrt{x}})
\end{equation}